\documentclass[11pt]{article}
\usepackage[margin=0.8in]{geometry}
\usepackage{amsmath,amssymb,amsthm}
\usepackage{enumitem}
\usepackage{longtable}
\usepackage{array}
\usepackage{xcolor}
\usepackage{hyperref}
\usepackage{fancyhdr}
\usepackage{titlesec}
\usepackage{tcolorbox}
\usepackage{multicol}
\usepackage{booktabs}
\hypersetup{colorlinks=true, linkcolor=blue!50!black, urlcolor=blue!50!black}
\pagestyle{fancy}
\fancyhf{}
\rhead{\thepage}
\lhead{Engineering Math to Pure Math}
\setlength{\parskip}{0.55em}
\setlength{\parindent}{0pt}
\setlist[itemize]{topsep=0.2em,itemsep=0.2em}
\setlist[enumerate]{topsep=0.2em,itemsep=0.25em}
\newtcolorbox{focusbox}{colback=blue!3!white,colframe=blue!50!black,title=Focus}
\newtcolorbox{intuitionbox}{colback=green!3!white,colframe=green!40!black,title=Intuition}
\newtcolorbox{formalbox}{colback=gray!4!white,colframe=gray!70!black,title=Formal}
\newtcolorbox{walkbox}{colback=orange!5!white,colframe=orange!60!black,title=Walking question}
\newcommand{\N}{\mathbb{N}}
\newcommand{\Z}{\mathbb{Z}}
\newcommand{\Q}{\mathbb{Q}}
\newcommand{\R}{\mathbb{R}}
\newcommand{\C}{\mathbb{C}}
\newcommand{\eps}{\varepsilon}
\newcommand{\pow}{\mathcal{P}}
\theoremstyle{definition}
\newtheorem{definition}{Definition}
\newtheorem{example}{Example}
\newtheorem{exercise}{Exercise}
\theoremstyle{plain}
\newtheorem{theorem}{Theorem}
\newtheorem{lemma}{Lemma}
\theoremstyle{remark}
\newtheorem{remark}{Remark}

\title{\textbf{Proof Workbook}\\Exercises, Guided Walkthroughs, and Thinking Prompts}
\author{Engineering Math to Pure Math}
\date{May 2026}
\begin{document}
\maketitle
\tableofcontents
\newpage

\section{How to Work These Problems}

Do not only hunt the answer. The goal is to train your eye. For each exercise, write the claim, identify the definitions, choose a strategy, and leave a readable proof.

\begin{focusbox}
A good beginner proof often feels almost too slow. That is fine. Slow proofs become fast understanding later.
\end{focusbox}

\section{Set 1: Logic and Translation}

\begin{exercise}
Translate each sentence into symbols. Then write its negation in symbols and in plain English.
\begin{enumerate}
\item Every integer has a square that is nonnegative.
\item Some real number is larger than its square.
\item For every programmer there is a bug that they cannot reproduce.
\item Every continuous function on $[0,1]$ is bounded.
\item There is a prime number larger than $1000$.
\end{enumerate}
\end{exercise}

\begin{exercise}
Determine whether each statement is true or false. If false, give a counterexample.
\begin{enumerate}
\item $\forall x\in\R,\exists y\in\R$ such that $x+y=0$.
\item $\exists y\in\R,\forall x\in\R$ such that $x+y=0$.
\item $\forall n\in\Z,\ n^2+n$ is even.
\item $\forall a,b\in\Z,\ a\mid b\Rightarrow b\mid a$.
\end{enumerate}
\end{exercise}

\section{Set 2: Direct Proof}

\begin{exercise}
Prove that if $a,b,c\in\Z$, $a\mid b$, and $b\mid c$, then $a\mid c$.
\end{exercise}

\textbf{Walkthrough.} Start from definitions. $a\mid b$ means $b=ak$ for some $k\in\Z$. $b\mid c$ means $c=b\ell$ for some $\ell\in\Z$. Substitute: $c=(ak)\ell=a(k\ell)$. Since $k\ell\in\Z$, $a\mid c$.

\begin{exercise}
Prove that the sum of two even integers is even.
\end{exercise}

\begin{exercise}
Prove that if $m,n\in\Z$ are odd, then $mn$ is odd.
\end{exercise}

\begin{exercise}
Prove that if $x\in\R$ and $0<x<1$, then $x^2<x$.
\end{exercise}

\section{Set 3: Contrapositive and Contradiction}

\begin{exercise}
Prove by contrapositive: if $n^2$ is even, then $n$ is even.
\end{exercise}

\begin{exercise}
Prove that $\sqrt{2}$ is irrational.
\end{exercise}

\begin{exercise}
Prove that there are infinitely many prime numbers.
\end{exercise}

\begin{walkbox}
For Euclid's proof of infinitely many primes, imagine the known primes as a fenced garden. What does multiplying them all and adding $1$ do to the fence?
\end{walkbox}

\section{Set 4: Sets}

\begin{exercise}
Prove $A\cap(B\cup C)=(A\cap B)\cup(A\cap C)$.
\end{exercise}

\textbf{Proof skeleton.}
\begin{enumerate}
\item To prove equality, prove both subset directions.
\item For $\subseteq$, let $x\in A\cap(B\cup C)$.
\item Then $x\in A$ and $x\in B\cup C$.
\item Split into cases: $x\in B$ or $x\in C$.
\item Finish the first direction, then reverse the reasoning for $\supseteq$.
\end{enumerate}

\begin{exercise}
Prove $A\setminus(B\cup C)=(A\setminus B)\cap(A\setminus C)$.
\end{exercise}

\begin{exercise}
Give an example where $A\subsetneq B$ and $\pow(A)\subsetneq\pow(B)$.
\end{exercise}

\section{Set 5: Functions}

\begin{exercise}
Let $f:A\to B$ and $g:B\to C$. Prove that if $f$ and $g$ are injective, then $g\circ f$ is injective.
\end{exercise}

\begin{exercise}
Let $f:A\to B$ and $S,T\subseteq B$. Prove
\[
f^{-1}(S\cap T)=f^{-1}(S)\cap f^{-1}(T).
\]
\end{exercise}

\begin{exercise}
Find a function $f:\R\to\R$ and sets $S,T\subseteq\R$ such that
\[
f(S\cap T)\neq f(S)\cap f(T).
\]
\end{exercise}

\section{Set 6: Relations}

\begin{exercise}
Define $a\sim b$ on $\Z$ by $a\sim b$ if $a-b$ is divisible by $5$. Prove $\sim$ is an equivalence relation.
\end{exercise}

\begin{exercise}
Let $R$ be a relation on a set $A$. Write the definitions of reflexive, symmetric, antisymmetric, and transitive. Give one relation that has each property and one that fails each property.
\end{exercise}

\section{Set 7: Induction}

\begin{exercise}
Prove that for every $n\in\N$,
\[
1+2+\cdots+n=\frac{n(n+1)}{2}.
\]
\end{exercise}

\begin{exercise}
Prove that for every $n\geq 1$,
\[
2^n\geq n+1.
\]
\end{exercise}

\begin{exercise}
Use strong induction to prove every integer $n\geq 2$ is a product of primes.
\end{exercise}

\section{Set 8: Countability}

\begin{exercise}
Prove that the set of even integers is countably infinite.
\end{exercise}

\begin{exercise}
Prove that $\Z$ is countable.
\end{exercise}

\begin{exercise}
Prove that $\Q$ is countable. Hint: arrange rational numbers by numerator-denominator grid and skip repeats.
\end{exercise}

\begin{exercise}
Explain Cantor's diagonal argument in your own words. Do not use the phrase ``it is obvious.''
\end{exercise}

\section{Set 9: First Analysis Problems}

\begin{exercise}
Prove directly from the definition that $\lim_{n\to\infty}1/n=0$.
\end{exercise}

\begin{exercise}
Prove that if $a_n\to a$ and $b_n\to b$, then $a_n+b_n\to a+b$.
\end{exercise}

\begin{exercise}
Prove that $f(x)=3x+2$ is continuous at every $c\in\R$ using the $\eps$-$\delta$ definition.
\end{exercise}

\section{Set 10: First Algebra Problems}

\begin{exercise}
Prove that the identity element of a group is unique.
\end{exercise}

\begin{exercise}
Prove that every element of a group has a unique inverse.
\end{exercise}

\begin{exercise}
Let $G$ be a group and $a,b\in G$. Prove $(ab)^{-1}=b^{-1}a^{-1}$.
\end{exercise}

\section{Portfolio Assignment}

Choose twelve exercises from this workbook: two logic, two set/function, two relation/induction, two countability/analysis, two algebra, and two wild cards. Rewrite them as a polished proof portfolio. For each proof, include:

\begin{itemize}
\item the claim,
\item the definitions used,
\item the proof strategy,
\item the final proof,
\item one sentence about the main idea.
\end{itemize}

\end{document}
